(1) 对半径 $r$ 的同轴圆周，$\oint\mathbf B\cdot d\boldsymbol\ell=\mu_0\varepsilon_0\,d\Phi_E/dt$。因而 $r<a$ 时 $B_\varphi=\mu_0\varepsilon_0r\dot E/2$，$r>a$ 时 $B_\varphi=\mu_0\varepsilon_0a^2\dot E/(2r)$；两式在 $r=a$ 连续，且环向取向由右手定则确定。
(2) 坡印廷矢量 $\mathbf S=\mathbf E\times\mathbf B/\mu_0$。因为 $\hat{\mathbf z}\times\hat{\boldsymbol\varphi}=-\hat{\mathbf r}$，边缘径向分量 $S_r(a)=-\varepsilon_0aE\dot E/2$，负号代表电磁能从侧面流入电容板间而非向外辐射。
(3) 侧面积为 $2\pi ad$，故总流入功率 $P_{\rm in}=-S_r(a)2\pi ad=\varepsilon_0\pi a^2d\,E\dot E$。板间电场能量 $U_E=\frac12\varepsilon_0E^2\pi a^2d$，微分即 $dU_E/dt=\varepsilon_0\pi a^2d\,E\dot E=P_{\rm in}$，验证局部电磁能守恒。
(4) 同一闭合安培回路可以跨过导线所在截面或板间截面，前者计传导电流，后者计位移电流，二者是同一充电过程在不同选定曲面上的表述，不可重复计数。准静态近似要求 $d\ll a$ 且充电时间尺度远长于 $a/c$；否则磁场和能流必须用含传播延迟的全麦克斯韦解。
